\textbf{COLUNI 2020 -- Questão 08}\par\medskip As personagens podem ser caracterizadas não só por suas descrições físicas e psicológicas, mas também pelo que falam e como falam, isto é, por seu discurso. Observe que, no fragmento abaixo, a fala da personagem apresenta-se na forma de discurso indireto livre.

\par

“Ele retirou a tampa com uma pinça, eu e meu irmão de olho, fungando curiosidade, e, \textit{Nossa!}, eis aquelas peças todas [...]” CARRASCOZA, J. A. \textit{Aos 7 e aos 40}. Rio de Janeiro: Alfaguara, 2016. p. 9.

\par

Assinale a alternativa em que esse mesmo recurso linguístico – discurso indireto livre – é utilizado da mesma forma nos trechos destacados:\begin{enumerate}
\item “[...] tia Maria preparava [...] ia me fazendo perguntas, \underline{Você tá na escola municipal ou na estadual?} [...]”
     (CARRASCOZA, J. A. \textit{Aos 7 e aos 40}. Rio de Janeiro: Alfaguara, 2016. p. 93.)
\item “Eu voava pra cozinha, entregue inteiramente à minha fome, \underline{Cuidado, fique longe da frigideira!}”
     (CARRASCOZA, J. A. \textit{Aos 7 e aos 40}. Rio de Janeiro: Alfaguara, 2016. p. 08.)
\item “Aquele dia de noite, o pai fazendo serão, / ela falou comigo: / ‘\underline{coitado, até essa hora no serviço pesado}’.”
     (PRADO, A. \textit{Bagagem}. 29ª ed. Rio de Janeiro. Record. 2019. p.118.)
\item “Um jogador de futebol chegou mesmo a declarar: / ‘\underline{Tenho birra de que me chamem de intelectual, / sou
      um homem como todos os outros}’.” (PRADO, A. \textit{Bagagem}. 29ª ed. Rio de Janeiro. Record. 2019. p. 40.)\end{enumerate}
